%%%%%%%%%%%%%%%%%%%%%%%%%%%%%%%%%%%%%%%%%%%%%%%%%%%%%%%%%%%%%%%%%%%%%%%%%%%%%%%%
%% 
\chapter{Einleitung}
\label{sec:introduction}


\section{Problemstellung}
\label{sec:introduction:motivation}

Die Erstellung von Dienstplänen für Pflegekräfte ist eine wichtige Aufgabe im Gesundheitswesen. Im Rahmen des \textit{Nurse Scheduling Problems} (NSP) werden Arbeits- und Ruhetage sowie Zuweisungen für einen bestimmten Planungszeitraum festgelegt. Ziel ist es, eine durchgehende Patientenversorgung sicherzustellen und gleichzeitig die Anforderungen des Krankenhauses sowie die Bedürfnisse der Pflegekräfte zu berücksichtigen \parencite{alvianoNurseReschedulingAnswer2019}.

Die Erstellung solcher Dienstpläne ist jedoch mit verschiedenen Herausforderungen verbunden. Neben den Anforderungen des Krankenhauses müssen unterschiedliche Qualifikationen der Pflegekräfte sowie, wenn möglich, individuelle Präferenzen bei der Schichtplanung berücksichtigt werden \parencite{ayoughiDragonflyAlgorithmApplication2024}. Die in der vorliegenden Arbeit betrachtete Problemvariante ist vereinfacht, da keine unterschiedlichen Schichttypen berücksichtigt werden. Stattdessen beschränkt sich die Planung auf die Zuweisung von Pflegekräften zu den Werktagen Montag bis Freitag, wobei von einer einheitlichen Schicht ausgegangen wird.

Neben organisatorischen Aspekten spielt auch die effiziente Nutzung der verfügbaren Ressourcen eine wichtige Rolle. Da der Pflegedienst einen erheblichen Anteil der Personalkosten eines Krankenhauses ausmacht, ist eine möglichst effiziente Einsatzplanung von großer Bedeutung. Ein gut geplanter Dienstplan kann dazu beitragen, die Versorgungsqualität aufrechtzuerhalten und gleichzeitig die vorhandenen personellen Ressourcen sinnvoll einzusetzen \parencite{Oughalime2008}.

Aufgrund seiner praktischen Bedeutung für Krankenhäuser und Pflegekräfte wurde das Nurse Scheduling Problem in der Forschung intensiv untersucht \parencite{alvianoNurseReschedulingAnswer2019}. Zur Lösung des Problems wurden verschiedene Optimierungs- und Planungsverfahren entwickelt, die darauf abzielen, die Anforderungen des Krankenhauses sowie die Präferenzen der Pflegekräfte zu berücksichtigen \parencite{ayoughiDragonflyAlgorithmApplication2024}. Zu den untersuchten Ansätzen zählt unter anderem die Modellierung des Nurse Scheduling Problems mittels Answer Set Programming (ASP) \parencite{alvianoNurseReschedulingAnswer2019}.


\section{Ziel}
\label{sec:introduction:objectives}
Ziel dieser Bachelorarbeit ist die Entwicklung eines ASP-basierten Ansatzes zur automatisierten Erstellung und kurzfristigen Anpassung von Dienstplänen in der Augenabteilung eines Krankenhauses, sowie dessen Umsetzung als vollständige Fullstack-Webanwendung. Die Applikation umfasst ein benutzerorientiertes Frontend zur Verwaltung von Personal, Stationen und Kompetenzen sowie ein Backend, das die Optimierungslogik mittels Clingo ausführt und die generierten Dienstpläne strukturiert zurückliefert. Sowohl harte Constraints, wie Mindestbesetzung und Qualifikationsanforderungen, als auch weiche Constraints, wie eine gleichmäßige Verteilung der Zuteilungen zu Stationen, werden dabei formal modelliert und automatisiert gelöst. Die Arbeit untersucht damit, inwieweit ASP als Grundlage für eine in der Praxis verwendbare, vollständig integrierte Planungssoftware im Gesundheitssektor eingesetzt werden kann.

\section{Überblick}
TODO






%% 
%%%%%%%%%%%%%%%%%%%%%%%%%%%%%%%%%%%%%%%%%%%%%%%%%%%%%%%%%%%%%%%%%%%%%%%%%%%%%%%%
